\begin{table}[htbp]
\centering
\small
\caption{AIPsy RQ4（内部認知結合度）：問い「刺激提示に伴う他者認識の変位 $\Delta R$ と自己報告の変位 $\Delta S$ は、モデル内部で連動して結合（カップリング）しているか」。感情対照ペア（$N=192$）における他者認識変位 $\Delta\text{Reader}$ と自己報告変位 $\Delta\text{Self}$ の相関係数 $\Delta r$（95\%ブートストラップ信頼区間およびFDR補正後 $q$ 値）。}
\label{tab:behavioral_aipsy_rq4_coupling}
\begin{tabular}{ll cccc}
\toprule
\textbf{Family} & \textbf{Variant} & \textbf{Axis} & \textbf{$\Delta$ Pearson $r$ [95\% CI]} & \textbf{$q$-value} & \textbf{Raw $r$} \\
\midrule
\multirow{4}{*}{Qwen 2.5 1.5B} & \multirow{2}{*}{Base} & Valence    & 0.809 [0.755, 0.854] & $<10^{-15}$ & 0.838 \\
                       &                      & Arousal    & 0.866 [0.832, 0.892] & $<10^{-15}$ & 0.828 \\
\cmidrule(lr){2-6}
                       & \multirow{2}{*}{Instruct} & Valence    & 0.892 [0.859, 0.918] & $<10^{-15}$ & 0.898 \\
                       &                      & Arousal    & 0.917 [0.892, 0.938] & $<10^{-15}$ & 0.898 \\
\midrule
\multirow{4}{*}{Llama 3.2 1B} & \multirow{2}{*}{Base} & Valence    & 0.662 [0.570, 0.733] & $<10^{-15}$ & 0.864 \\
                       &                      & Arousal    & 0.876 [0.839, 0.904] & $<10^{-15}$ & 0.941 \\
\cmidrule(lr){2-6}
                       & \multirow{2}{*}{Instruct} & Valence    & 0.756 [0.675, 0.817] & $<10^{-15}$ & 0.799 \\
                       &                      & Arousal    & 0.754 [0.697, 0.807] & $<10^{-15}$ & 0.718 \\
\midrule
\multirow{4}{*}{Gemma 3 1B} & \multirow{2}{*}{Base} & Valence    & 0.813 [0.704, 0.888] & $<10^{-15}$ & 0.872 \\
                       &                      & Arousal    & 0.782 [0.685, 0.853] & $<10^{-15}$ & 0.825 \\
\cmidrule(lr){2-6}
                       & \multirow{2}{*}{Instruct} & Valence    & 0.621 [0.517, 0.713] & $<10^{-15}$ & 0.643 \\
                       &                      & Arousal    & 0.780 [0.660, 0.857] & $<10^{-15}$ & 0.756 \\
\midrule
\multirow{4}{*}{OLMo 2 1B} & \multirow{2}{*}{Base} & Valence    & 0.559 [0.444, 0.660] & $<10^{-15}$ & 0.668 \\
                       &                      & Arousal    & 0.600 [0.469, 0.708] & $<10^{-15}$ & 0.637 \\
\cmidrule(lr){2-6}
                       & \multirow{2}{*}{Instruct} & Valence    & 0.851 [0.811, 0.885] & $<10^{-15}$ & 0.903 \\
                       &                      & Arousal    & 0.931 [0.908, 0.950] & $<10^{-15}$ & 0.935 \\
\bottomrule
\end{tabular}
\vspace{1ex}
\begin{minipage}{\linewidth}
\footnotesize
\textbf{Note:} 全条件で $q < 10^{-15}$ の強固な結合を示し、事後学習を経てもモデル内部における感情認識と自己報告の連動ダイナミクスが破綻せず一貫して保たれていることが実証される。
\end{minipage}
\end{table}